\originalsection{作业11: 2011年12月8日交}
\begin{exercise}\label{ps:11I}
\textbf{原题 I.} 可微可逆矩阵 $N(q)$, 证明 $(N^{-1})'=-N^{-1}N'N^{-1}$.
\end{exercise}
\begin{solution}
微分 $N^{-1}N=I$ 得 $(N^{-1})'N+N^{-1}N'=0$. 右乘 $N^{-1}$ 即结论; 按分量写为
\[
\partial_q(N^{-1})^{\mu\nu}=-(N^{-1})^{\mu\alpha}(\partial_qN_{\alpha\beta})(N^{-1})^{\beta\nu}.
\]
\end{solution}
\begin{exercise}\label{ps:11II}
\textbf{原题 II(a--d).} 用 Frobenius 范数 $\|A\|=(\sum_{ij}|A_{ij}|^2)^{1/2}$, 证明次乘性; 小矩阵的行列式展开; 有限几何级数; 充分小时的 Neumann 逆级数.
\end{exercise}
\begin{solution}
(a)逐项 Cauchy--Schwarz 给 $|(AB)_{ij}|^2\le(\sum_k|A_{ik}|^2)(\sum_k|B_{kj}|^2)$, 对i,j求和得 $\|AB\|^2\le\|A\|^2\|B\|^2$.
(b)Leibniz 行列式公式中常数项为1, 一次项仅对角 $\sum_iA_{ii}=\operatorname{tr}A$, 其余有限多单项式次数至少2. 因 $|A_{ij}|\le\|A\|$, 在 $\|A\|\le1$ 有
$\det(I+A)=1+\operatorname{tr}A+\|A\|^2R(A)$, $R$ 有界. 原卷首项写矩阵I, 应为标量1.
(c)直接展开相消得 $(I-A)\sum_{k=0}^NA^k=I-A^{N+1}$, 因A的幂互相交换, 左右次序均可.
(d)由(b)及 $|\operatorname{tr}A|\le\sqrt{n+1}\|A\|$ 可选 $\gamma<1$ 使 $\det(I-A)\ne0$. 又(a)给 $\|A^k\|\le\gamma^k$ ($k\ge1$), 逆级数绝对收敛于矩阵S; (c)取极限给 $(I-A)S=S(I-A)=I$, 所以 $S=(I-A)^{-1}$. 实际单用 $\|A\|<1$ 的级数已能直接证明可逆.
\end{solution}
\begin{exercise}\label{ps:11III}
\textbf{原题 III(a--d).} $\mathcal L=-m^{\alpha\beta}\phi_\alpha\phi_\beta/2-V(\phi)$, $V$ 光滑实值. 说明坐标协变, 推场方程, 求应力张量并直接证明无散度.
\end{exercise}
\begin{solution}
(a) $\phi$ 为标量, $\phi_\alpha$ 为协向量, 度量逆按双逆指标变换, 收缩 $m^{\alpha\beta}\phi_\alpha\phi_\beta$ 为标量, $V(\phi)$ 亦然. 任意坐标下度量分量须同步变换; 不能变坐标却仍强行令矩阵diag不变. 不变作用量应配体积密度, 在 Lorentz 坐标中Jacobian绝对值1.
(b)紧支撑变分h给 $\delta S=\int[-\partial^\alpha\phi\partial_\alpha h-V'(\phi)h]=\int(\Box\phi-V'(\phi))h$. 故 $\Box\phi=V'(\phi)$, 即 $\phi_{tt}-\Delta\phi+V'(\phi)=0$.
(c)度量变分或第十二章张量公式给
$T^{\mu\nu}=\partial^\mu\phi\partial^\nu\phi-m^{\mu\nu}(Q/2+V(\phi))$.
(d)乘积求导后二阶交叉项与 $\partial^\nu Q/2$ 抵消,
$\partial_\mu T^{\mu\nu}=(\Box\phi-V'(\phi))\partial^\nu\phi=0$.
若 $V=\kappa^2\phi^2/2$, 场方程为质量参数 $\kappa$ 的 Klein--Gordon 方程.
\end{solution}
\begin{exercise}\label{ps:11IV}
\textbf{原题 IV(a--d), 含勘误.} 设 $\phi$ 为原题 III 场方程的 $C^2$ 解, $T$ 为其应力张量. $Y$ 的流 $\widetilde x$ 保持 Minkowski 度量, $M=\partial\widetilde x/\partial x$. (a)微分 $Mm^{-1}M^T=m^{-1}$ 得 Killing 方程; (b)证 $\partial_\mu(T^{\mu\nu}Y_\nu)=0$ 及其区域积分式; (c)解释守恒量; (d)取 $Y=(-1,0,\ldots)$ 得带势能量守恒.
\end{exercise}
\begin{solution}
(a)在流参数0处 $M=I$, $\partial_\eps M^\mu_\alpha=\partial_\alpha Y^\mu$. 微分后
$m^{\alpha\nu}\partial_\alpha Y^\mu+m^{\mu\alpha}\partial_\alpha Y^\nu=0$.
降指标得 $\partial_\mu Y_\nu+\partial_\nu Y_\mu=0$.
(b)使用上题无散度及对称性,
$\partial_\mu(T^{\mu\nu}Y_\nu)=T^{\mu\nu}\partial_\mu Y_\nu=\tfrac12T^{\mu\nu}(\partial_\mu Y_\nu+\partial_\nu Y_\mu)=0$.
对任何积分存在的有界光滑区域K积分也为0. 原卷用“任意光滑Z”的变分积分恒等式, 实际须Z内部紧支撑或有适当边界条件; 这里直接证目标, 不沿用错误的无边界限制变分.
(c)对 $[t_1,t_2]\times B_R$ 积分散度, 得两时间片积分差等于负侧面通量. 若侧面通量为0或 $R\to\infty$ 时趋0, 则 $\int T^{0\nu}Y_\nu$ 守恒. 这是等距对称产生守恒量的具体形式.
(d) $Y_0=1$, 故流为 $T^{\mu0}$, 时间分量 $T^{00}=(\phi_t^2+|\nabla\phi|^2)/2+V(\phi)$. 取初值紧支撑、有限时间段支撑受控的解并假设 $V(0)=V'(0)=0$, 得
\[
\int\left[\tfrac12(\phi_t^2+|\nabla\phi|^2)+V(\phi)\right](t,x)\dd x
=\int\left[\tfrac12(\phi_t^2+|\nabla\phi|^2)+V(\phi)\right](0,x)\dd x.
\]
原卷右边梯度误留t, 此处全部在0评价; 场方程应回指原(0.0.8), 而非 Lagrangian(0.0.7). 若 $V(0)\ne0$, 必须用 $V(\phi)-V(0)$ 消除无限常数势能; 无穷远条件亦不可省.
\end{solution}
